%% Compile: pdflatex paper_GP_MFPCA  (twice for cross-references)
\documentclass[12pt,a4paper]{article}

\usepackage{mathptmx}
\usepackage[T1]{fontenc}
\usepackage[utf8]{inputenc}
\usepackage{geometry}
\geometry{a4paper, top=25mm, bottom=25mm, left=25mm, right=25mm}
\usepackage{amsmath,amssymb}
\usepackage{graphicx}
\usepackage{booktabs}
\usepackage{array}
\usepackage{xcolor}
\usepackage{tikz}
\usetikzlibrary{shapes.geometric, arrows.meta, positioning, fit, backgrounds}
\usepackage{natbib}
\usepackage{setspace}
\usepackage{caption}
\usepackage{hyperref}
\hypersetup{colorlinks=true, linkcolor=blue!60!black,
            citecolor=blue!60!black, urlcolor=blue}
\usepackage{lineno}
\usepackage{authblk}
\usepackage{titlesec}
\usepackage{parskip}

\setstretch{1.15}

\titleformat{\section}{\large\bfseries}{\thesection.}{0.5em}{}
\titleformat{\subsection}{\normalsize\bfseries}{\thesubsection.}{0.5em}{}
\titleformat{\subsubsection}{\normalsize\itshape}{\thesubsubsection.}{0.5em}{}

%% ------------------------------------------------------------------
%% Notation:
%%   \spar        -- theta == Sigma_{z,1}: initial solid surface density
%%                  at 1 au (g cm^{-2}); primary simulation control parameter.
%%                  Range [23.39, 175.88] g cm^{-2}. Defined in Sec. 2.1.
%%   v_k          -- GP interpolation posterior variance for the smooth score
%%   \delta_k^2   -- fixed per-PC LOESS-SE^2 observation noise used in the GP
%%   \sigma^2_{\mathrm{resid},k} -- per-PC LOESS residual variance restored
%%                  at prediction time
%%   \sigma^2_{\mathrm{pred},k} -- total predictive score variance
%%                  v_k + \sigma^2_{\mathrm{resid},k}
%% ------------------------------------------------------------------
\newcommand{\GP}{\textsc{gp}}
\newcommand{\MFPCA}{\textsc{mfpca}}
\newcommand{\KDE}{\textsc{kde}}
\newcommand{\CLR}{\textsc{clr}}
\newcommand{\SQRT}{\ensuremath{\sqrt{f}}}
\newcommand{\ISE}{\textsc{ise}}
\newcommand{\CRPS}{\textsc{crps}}
\newcommand{\iCRPS}{i\textsc{crps}}
\newcommand{\ECDF}{\textsc{ecdf}}
\newcommand{\calK}{\mathcal{K}}
\newcommand{\spar}{\theta}

%% ================================================================
\title{\textbf{Emulating Multivariate Conditional Densities from
$N$-body Simulations: A Gaussian Process Pipeline with Calibrated
Uncertainty Quantification}}

\author[1]{Yiming Chen}
\affil[1]{Independent researcher, United States}

\date{\small\today}

%% ================================================================
\begin{document}
\linenumbers
\maketitle
\thispagestyle{empty}

%% ----------------------------------------------------------------
\begin{abstract}
\noindent
We design and validate a Gaussian process emulation pipeline for
multivariate conditional density functions from $N$-body planetary
system simulations.  The emulator serves as a fast surrogate for
Bayesian inference of planet formation conditions from Kepler
observations; the present paper establishes its predictive performance
against held-out simulations.

The pipeline combines adaptive kernel density estimation, a
square-root density (Hellinger) transform, multivariate functional
PCA, LOESS denoising, and GP regression with ECDF input warping and a
noise-decomposition uncertainty model.  A key architectural insight is
the choice of density transform: replacing the standard centred
log-ratio (CLR) transform with the square-root transform eliminates the
exponential error amplification inherent in the CLR inverse.  The
resulting two-component predictive variance---GP interpolation
uncertainty plus restored LOESS residual variance---achieves
per-variable density-space coverage of 94.1\,\% (log period ratio),
84.9\,\% (log duration ratio), and 94.0\,\% (log Hill spacing) at
95\,\% nominal level, without post-hoc calibration.

Applied to three planetary-system observables conditioned on the solid
surface density parameter $\spar$, the emulator achieves statistically
significant improvements on the integrated CRPS (\iCRPS) proper scoring
rule for all three observables ($p<0.05$, $n=25$), with ratios of 0.84
(log period ratio), 0.76 (log duration ratio), and 0.81 (log Hill
spacing) relative to the training-mean naive baseline.  An error decomposition
identifies GP score prediction (not basis representation) as the
dominant remaining error source, accounting for 60--74\,\% of total
density-space ISE.
\end{abstract}

\textbf{Keywords:} methods: statistical --- methods: numerical ---
planets and satellites: dynamical evolution and stability ---
methods: data analysis --- Gaussian processes --- density estimation

\vspace{0.5cm}\hrule\vspace{0.5cm}

%% ================================================================
\section{Introduction}
\label{sec:intro}

Population-level statistics of planetary systems encode the cumulative
outcome of formation and dynamical-evolution processes operating over
millions to billions of years \citep{winn2015, mordasini2018,
mulders2021}.  A central scientific goal is to \emph{infer} the
conditions under which a planetary system formed---disc surface density,
radial distribution of solids, gas disk depletion---from the observable
properties of that system (orbital periods, transit durations, Hill
spacings) recorded by missions such as \textit{Kepler} and \textit{K2}.

This inverse problem requires repeated evaluation of the likelihood
$P(\text{obs}\,|\,\boldsymbol{\phi})$, which in turn requires evaluating the
forward model $F(\boldsymbol{\phi})$ that maps formation parameters
$\boldsymbol{\phi}=(\theta,\gamma,d)$ to a distribution of observables.
Here $\theta\equiv\Sigma_{z,1}$ is the solid surface density at 1\,au,
$\gamma$ is the disc radial exponent (held fixed at $\gamma=-3/2$ in this paper), and $d$ is the gas disc depletion factor;
the present paper holds $\gamma$ and $d$ fixed and treats $\theta$ as the
sole free parameter (Section~\ref{sec:simsetup}).  The forward model here is a full $N$-body
gravitational simulation: each run requires approximately five days of
wall-clock time \citep{macdonald2020}, making MCMC-based posterior
inference---which requires thousands of likelihood evaluations---
computationally infeasible \citep{ida2004, benz2014, emsenhuber2021}.

The standard remedy is to train a statistical \emph{emulator}
\citep{kennedy2001, bower2010} on an existing design of simulation runs,
producing a fast surrogate that can be queried millions of times during
MCMC at negligible cost.  The present paper designs and validates such an
emulator.  Application to Bayesian inference of planet formation
conditions from Kepler data is the subject of ongoing work; the emulator
must be validated before it can be trusted in the inverse-problem
context.

The emulation task here is non-standard: the simulator output at each
parameter setting $\spar$ is not a scalar but a
\emph{multivariate probability density function}---the joint
conditional distribution of three observable planetary-system
properties (orbital spacing, transit-duration ratio, and dynamical
proximity) given $\spar$.  Standard scalar-output Gaussian process
emulators \citep{rasmussen2006} do not handle functional outputs
directly.  Latent-space methods compress the functional output into a
low-dimensional score vector via functional PCA and emulate the scores
\citep{higdon2008, fricker2013}; this approach is adopted here,
extended to the multivariate functional setting via \MFPCA
\citep{happ2018}.

\subsection*{Contribution}

This paper designs and \emph{validates} a Gaussian process emulation
pipeline for multivariate conditional density functions from $N$-body
planetary system simulations.  The pipeline combines adaptive \KDE,
the square-root density transform, \MFPCA, and a \GP with \ECDF input
warping into a fully specified, reproducible workflow (all
hyperparameters in Table~\ref{tab:hyperparams}).  The contribution has
three parts:
\begin{enumerate}
  \item \emph{Pipeline design with near-calibrated uncertainty propagation.}
        The pipeline achieves per-variable empirical coverage of
        94.1\,\% ($X$), 84.9\,\% ($Y$), and 94.0\,\% ($Z$) at
        95\,\% nominal level in density space, without post-hoc
        calibration.
        We identify the square-root density transform as the critical
        architectural choice enabling this calibration, and justify the
        homoscedastic GP with fixed LOESS-SE$^2$ observation noise and
        explicit residual-variance restoration as the appropriate noise
        model for the two-stage (LOESS denoise $\to$ GP emulate)
        architecture.
  \item \emph{Multi-baseline empirical evaluation.}  We test the
        pipeline against a parameter-ignorant naive baseline, a
        score-space LOESS interpolant, and a $\spar$-aware
        nearest-neighbour estimator, finding statistically significant
        \iCRPS gains for all three observables ($p<0.05$;
        Section~\ref{sec:density_eval}).
  \item \emph{Ablation analysis.}  We ablate three design choices---
        density transform (CLR vs $\sqrt{f}$), peak alignment (with vs
        without), and GP noise model (the final noise-decomposition
        model vs no-warp and heteroscedastic alternatives)---identifying
        the transform choice as the dominant factor in density-space
        calibration (coverage: 61\,\% with CLR vs near-nominal coverage
        for $X$ and $Z$ with $\sqrt{f}$).
\end{enumerate}

The remainder of the paper validates this emulator rather than
applying it to data inference.  The inverse problem---using the
validated emulator as a likelihood surrogate within MCMC to infer
$\boldsymbol{\phi}=(\theta,\gamma,d)$ from Kepler observables---is reserved for subsequent work.
Section~\ref{sec:data} describes the simulated data.
Section~\ref{sec:methods} details the pipeline.
Section~\ref{sec:results} presents results.
Section~\ref{sec:discussion} interprets the results and characterises
sensitivity.  Section~\ref{sec:conclusions} concludes.

%% ================================================================
\section{Data and Simulations}
\label{sec:data}

\subsection{Simulation setup}
\label{sec:simsetup}

The dataset comprises $N$-body planetary system simulations
parameterised by a single scalar control parameter $\theta$, the
initial solid surface density at 1\,au,
$\theta\equiv\Sigma_{z,1}$ (g\,cm$^{-2}$).
This normalises the total amount of solid material available
for late-stage planet formation in the inner disc.  Specifically,
the solid surface density profile follows
\begin{equation}
  \Sigma(a) = \theta\left(\frac{a}{\mathrm{au}}\right)^{-3/2},
  \label{eq:sigma_profile}
\end{equation}
adopting a minimum-mass solar nebula radial exponent $\gamma=-3/2$ (held fixed in this paper).
All other simulation parameters are held fixed: the gas disc is in
a strongly depleted state ($d=100$ in the gas surface density
prescription $\Sigma_g=1700\,\mathrm{g\,cm^{-2}}\,d^{-1}
(a/\mathrm{au})^{-3/2}$, using a step-function dispersal
approximation), and the initial embryo masses and spacings correspond
to 3~Hill radii.  Larger~$\spar$ implies more massive planetary
embryos and, consequently, more compact and massive post-formation
systems; the parameter therefore acts as the primary driver of
structural diversity in the simulated planetary population
\citep{macdonald2020}.

Simulations span $\spar\in[23.39,175.88]\,\mathrm{g\,cm^{-2}}$.
The canonical minimum-mass solar nebula (MMSN) corresponds to
$\theta_\mathrm{MMSN}\equiv\Sigma_{z,\mathrm{MMSN}}(1\,\mathrm{au})
\approx7\,\mathrm{g\,cm^{-2}}$ \citep{hayashi1981},
so the simulation range $\theta\in[23.39,175.88]\,\mathrm{g\,cm^{-2}}$
spans approximately 3--25 times the MMSN value.
After filtering levels
with fewer than 40 realisations to ensure reliable \KDE, 124 distinct
$\spar$ values remain, each with at least 40 independent system
outcomes.  These 124 values are partitioned at the $\spar$ level into
99 training and 25 held-out test values (approximately 80/20).  The
implicit smoothness assumption---that the three density functions vary
smoothly with $\spar$ on the scale of the inter-training spacing---is
verified empirically via the lag-1 autocorrelation of the smoothed PC
scores (Section~\ref{sec:sign}).

\subsection{Observables}
\label{sec:observables}

For each simulated system we extract three quantities for adjacent
planet pairs:
\begin{align}
  X &= \log\!\bigl(P_{i+1}/P_i\bigr), \label{eq:X}\\
  Y &= \log\!\bigl(\tau_{i+1}/\tau_i\bigr), \label{eq:Y}\\
  Z &= (a_{i+1}-a_i)/R_{\mathrm{H}}, \label{eq:Z}
\end{align}
where $P_i$ is orbital period, $\tau_i$ transit duration, $a_i$
semi-major axis, and $R_{\mathrm{H}}$ the mutual Hill radius.  These
are standard diagnostics in observational analyses of \textit{Kepler}
and \textit{K2} catalogues \citep{fabrycky2014, weiss2018, gilbert2020}.

%% ================================================================
\section{Pipeline}
\label{sec:methods}

\subsection{Inferential target and uncertainty decomposition}
\label{sec:target}

\paragraph{Inferential target.}
The emulator targets the \emph{conditional mean density}:
\begin{equation}
  f^*(t|\spar) = \lim_{n\to\infty}\hat{f}(t|\spar),
  \label{eq:target}
\end{equation}
the density the simulator produces at $\spar$ in the large-sample
limit, estimated via \KDE from a finite sample of $n(\spar)$
realisations.  The \GP posterior mean approximates this object.
Posterior predictive samples represent uncertainty \emph{about this
conditional mean density}---not individual realisation
variability---and should be interpreted accordingly.

\paragraph{Uncertainty decomposition.}
\textit{Note.} The following decomposition uses notation defined in
Sections~\ref{sec:mfpca}--\ref{sec:noise_decomp}; readers may find it useful
to return here after reading those sections.
Four sources of uncertainty arise; confounding them leads to both
over- and under-confident interval claims.
\begin{enumerate}
  \item \emph{GP interpolation uncertainty} ($v_k$; defined in
        Section~\ref{sec:gp}): posterior variance of the conditional
        mean \MFPCA score at an unobserved $\spar^*$.
  \item \emph{Residual score variability} ($s_k^2$; defined in
        Section~\ref{sec:noise_decomp}): variability of raw \MFPCA scores
        (Section~\ref{sec:mfpca}) around the LOESS-smoothed trend
        (Section~\ref{sec:gp})---a mixture of \KDE sampling noise
        and genuine within-$\spar$ distributional variation.  The two
        cannot be separated without replicated density estimates at
        the same~$\spar$.
  \item \emph{Basis representation error} (defined in
        Section~\ref{sec:mfpca}): the smooth B-spline basis cannot
        resolve arbitrarily narrow density features.  This is a
        systematic bias, not stochastic uncertainty; it is quantified
        by the \MFPCA reconstruction floor (Table~\ref{tab:mfpca_recon},
        Section~\ref{sec:mfpca_quality}).
  \item \emph{Noise-decomposition approximation error} (defined in
        Section~\ref{sec:noise_decomp}): the approximation
        $v_k(\spar^*)+\sigma^2_{\mathrm{resid},k}$ assumes that the
        LOESS residual variance adequately estimates KDE sampling noise
        and is approximately constant across $\spar$.  It also assumes
        approximate independence between GP interpolation error and KDE
        sampling noise.  These are simplifications; their empirical
        adequacy is assessed in
        Sections~\ref{sec:noise_coverage} and~\ref{sec:ablation_transform}.
\end{enumerate}
Components 1 and 2 are propagated via the noise-decomposition
predictive variance (Section~\ref{sec:noise_decomp}): the GP posterior variance
$v_k$ captures component~1, while the restored LOESS residual variance
$\sigma^2_{\mathrm{resid},k}$ captures component~2.
Component 3 sets the representation floor quantified in
Table~\ref{tab:mfpca_recon}.  Component 4 is controlled empirically:
with the square-root density transform, this yields density-space
coverage of 94.1\,\% ($X$), 84.9\,\% ($Y$), and 94.0\,\% ($Z$) at
95\,\% nominal level, without post-hoc calibration.

\subsection{Adaptive kernel density estimation}
\label{sec:kde}

At each $\spar$ level, $n(\spar)$ realisations are converted to smooth
density estimates on fixed 200-point grids spanning the empirical
training range.  All integrals are approximated via the trapezoidal
rule.

\paragraph{Pilot bandwidth (Silverman's rule \citep{silverman1986}).}
For each observable $v\in\{X,Y,Z\}$, the pilot bandwidth at training
level $\spar$ is
\begin{equation}
  h_0(\spar) = 0.9\min\!\bigl(s_v(\spar),\;
               {\mathrm{IQR}_v(\spar)}/{1.34}\bigr)
               \cdot n(\spar)^{-1/5},
  \label{eq:silverman}
\end{equation}
where $s_v(\spar)$ and $\mathrm{IQR}_v(\spar)$ are the sample standard
deviation and interquartile range of the $n(\spar)$ realisations of
observable $v$ at level $\spar$.

\paragraph{Abramson adaptive scaling \citep{abramson1982}.}
Here $x_j$ ($j=1,\ldots,n(\spar)$) denotes the $j$-th simulated
observation (e.g., log period ratio for one planet pair in one
realisation at the given $\spar$ level).
Let $\tilde{f}(\cdot|\spar)$ denote the pilot \KDE evaluated with
the uniform bandwidth $h_0(\spar)$.  The per-observation scale factor is
\begin{equation}
  \lambda_j =
  \bigl(\tilde{f}(x_j|\spar)/G_{\mathrm{geom}}\bigr)^{-1/2},\quad
  G_{\mathrm{geom}} = \exp\!\bigl(n^{-1}\textstyle\sum_j
                      \log\tilde{f}(x_j|\spar)\bigr).
  \label{eq:abramson}
\end{equation}
The final estimate with effective bandwidth $h_j=h_0\lambda_j$ is
\begin{equation}
  \hat{f}(t|\spar) = \frac{1}{n(\spar)}\sum_j
  \frac{1}{h_j}\kappa\!\bigl((t-x_j)/h_j\bigr),
  \label{eq:kde}
\end{equation}
where $\kappa$ is the Gaussian kernel.  A dynamic support mask zeroes and renormalises density at grid points outside the observed data range extended by 10\,\% of the empirical range on each side (i.e., beyond $[\min(x_j)-0.1\cdot R,\;\max(x_j)+0.1\cdot R]$ where $R=\max(x_j)-\min(x_j)$).
Grid values are then smoothed with an 80-knot B-spline basis
(penalty $\lambda=10^{-5}$) using the \texttt{fda} package
\citep{ramsay2005}.

\subsection{Square-root density (Hellinger) transform}
\label{sec:sqrt_transform}

Probability densities are compositional: non-negative and integrating
to unity.  Applying FPCA directly to raw densities can produce
reconstructions that violate non-negativity.  A standard remedy is
the centred log-ratio (CLR) transform, which maps densities to
unconstrained $L^2$ functions.  However, the CLR inverse involves
exponentiation, which \emph{exponentially amplifies} small
score-prediction errors at density peaks---precisely the regions where
accuracy matters most (Section~\ref{sec:ablation_transform}).

We instead apply the \emph{square-root density} (Hellinger) transform:
\begin{equation}
  h(t) = \sqrt{f(t)},
  \label{eq:sqrt_fwd}
\end{equation}
and work with $h(\cdot|\spar)$ throughout the MFPCA and GP stages.
The inverse is:
\begin{equation}
  f(t) = h(t)^2 \Big/ \int_a^b h(s)^2\,\mathrm{d}s,
  \label{eq:sqrt_inv}
\end{equation}
which guarantees non-negativity and involves only quadratic---not
exponential---error amplification.  A score-prediction error of
magnitude $\epsilon$ in the square-root space produces a density error
of order $\epsilon^2$ at the peak, compared to $e^\epsilon-1\approx
\epsilon$ under CLR, with a multiplicative constant that grows
exponentially with peak height.

Values below $\varepsilon_\mathrm{floor} = 5\times10^{-3}$ of the
local maximum are floored before the transform to prevent numerical
issues in the sparse tails.

The square-root transform operates in the Hellinger geometry
\citep{srivastava2007}, which is a natural metric for comparing
probability densities.  Unlike CLR, it does not map to a zero-integral
subspace; the MFPCA mean function and eigenfunctions are defined in
the square-root density space and need not integrate to zero.

\subsection{Multivariate functional PCA}
\label{sec:mfpca}

The square-root-transformed densities $h_X$, $h_Y$, $h_Z$ are
functions of different variables on different domains.  \MFPCA
\citep{happ2018} decomposes the joint variation of all three
square-root densities simultaneously:
\begin{equation}
  \mathbf{h}(\cdot|\spar)
  \approx\boldsymbol{\mu}(\cdot)
  +\sum_{k=1}^{K}\xi_k(\spar)\,\boldsymbol{\phi}_k(\cdot),
  \label{eq:mfpca}
\end{equation}
where $\boldsymbol{\mu}$ is the cross-$\spar$ mean function in
square-root density space, $\boldsymbol{\phi}_k$ are multivariate
eigenfunctions with equal variable weights ($w_v=1/3$), each
represented in a $n_\phi=50$ B-spline basis, and $\xi_k(\spar)$ are
scalar PC scores.  We retain $K=30$ components.

No curve registration is applied.  An earlier version of this pipeline
registered the dominant resonance location in $X$ and $Z$ before
MFPCA, but the geometry parameters required emulation at test $\spar$
values via LOESS, introducing a dominant error source (63--71\,\% of
density-space ISE).  Removing that registration step eliminates this
bottleneck entirely; the MFPCA basis absorbs resonance movement
directly through additional eigenfunctions.

\subsection{Sign alignment and LOESS pre-smoothing}
\label{sec:sign}

\paragraph{Sign alignment.}
Eigenfunctions are sign-indeterminate: $(\boldsymbol{\phi}_k,\xi_k)$
and $(-\boldsymbol{\phi}_k,-\xi_k)$ describe the same reconstruction.
We adopt the convention that the sign is chosen so that the score
series $\{\xi_k(\theta_j)\}$ has a non-negative Pearson correlation
with $\{\theta_j\}$; concretely, we negate both $\boldsymbol{\phi}_k$
and all $\xi_k$ whenever this correlation is negative.  The rationale
is purely presentational: scores that tend to increase with $\spar$
are easier to interpret visually (Figure~\ref{fig:scores}).  This
convention does not alter the reconstruction subspace, fitted densities,
or any metric value, and no monotonicity of $\xi_k$ in $\spar$ is
implied.

\paragraph{LOESS pre-smoothing.}
Raw scores $\xi_k(\spar)$ contain \KDE sampling noise which, if left
in the training signal, inflates GP hyperparameter estimates and
degrades mean-function accuracy.  We apply LOESS (span
$s_\mathrm{loess}=0.25$):
\begin{equation}
  \tilde{\xi}_k(\spar)=\mathrm{LOESS}_{0.25}[\xi_k](\spar).
  \label{eq:loess}
\end{equation}
This denoising step improves GP mean estimation but discards part of
the true within-$\spar$ variability.  The discarded variance
$\sigma^2_{\mathrm{resid},k}=\mathrm{Var}[r_k]$ is restored at
prediction time (Section~\ref{sec:noise_decomp}) as a fixed additive component
of the predictive variance.
The lag-1 autocorrelation of the smoothed series ranges from 0.92 to
0.98 across PCs~1--8
(Figure~\ref{fig:scores}), confirming structured $\spar$-dependence
throughout.

The raw-minus-smoothed residuals,
\begin{equation}
  r_k(\spar)=\xi_k(\spar)-\tilde{\xi}_k(\spar),
  \label{eq:residuals}
\end{equation}
are retained for two purposes: (i) estimating the GP observation noise
(Section~\ref{sec:gp}), and (ii) restoring the discarded variance at
prediction time (Section~\ref{sec:noise_decomp}).

\subsection{Homoscedastic Gaussian process with fixed known observation noise}
\label{sec:gp}

\paragraph{ECDF input warping.}
The 99 training $\spar$ values are mapped to approximately uniform
density in $[0,1]$ via the empirical CDF:
\begin{equation}
  u(\spar)=\frac{1}{99}\sum_{j=1}^{99}\mathbf{1}[\theta_j\le\spar],
  \label{eq:ecdf}
\end{equation}
This maps the $j$-th training point to $j/99$, giving a range $[1/99,1]$; the values are then linearly rescaled to $[0,1]$ before standardisation to zero mean and unit variance.

\paragraph{Two-stage smoother motivation.}
Raw MFPCA scores contain substantial KDE sampling noise: mean
noise SD across components is 0.44, while signal SD (the smooth
$\spar$-dependent trend) is of comparable magnitude, yielding
signal-to-noise ratios as low as 0.2 for leading PCs.  At this
noise level with $n=99$ training points, a GP cannot reliably
separate signal from noise; empirically, fitting a GP directly on raw
scores with the same fixed observation-noise rule produces flat predictions with
$\sigma^2\approx0$.  LOESS pre-smoothing is therefore a necessary
first-stage denoiser; the GP emulates the denoised curve with the
LOESS estimation uncertainty declared as fixed observation noise.

\paragraph{Model.}
For each component $k=1,\ldots,30$:
\begin{equation}
  \tilde{\xi}_k(\spar)\sim\GP\!\bigl(\beta_{0,k}+\beta_{1,k}u,\;
    \alpha_k^2\,\calK_{\ell_k}(u,u')+\delta_k^2\bigr),
  \label{eq:gp}
\end{equation}
where $\beta_{0,k}+\beta_{1,k}u$ is a linear mean function,
$\alpha_k^2$ the signal variance, $\calK_{\ell_k}$ the
Mat\'{e}rn-$\tfrac{3}{2}$ kernel:
\begin{equation}
  \calK_\ell(u,u')=
  \Bigl(1+\frac{\sqrt{3}|u-u'|}{\ell}\Bigr)
  \exp\!\Bigl(-\frac{\sqrt{3}|u-u'|}{\ell}\Bigr).
  \label{eq:matern32}
\end{equation}
The length-scale is bounded below at $\ell_\mathrm{min}=0.08$ in
\ECDF-warped space.

\paragraph{Fixed observation noise.}
The observation noise is fixed from the LOESS estimation uncertainty:
\begin{equation}
  \delta_k^2 =
  \frac{\sigma^2_{\mathrm{resid},k}}{n_\mathrm{eff}},
  \qquad n_\mathrm{eff} = \lfloor s_\mathrm{loess}\times N_\mathrm{train}\rfloor,
  \label{eq:obs_noise}
\end{equation}
where
$\sigma^2_{\mathrm{resid},k}=\mathrm{Var}[\xi_k-\tilde{\xi}_k]$
is the per-component LOESS residual variance and
$n_\mathrm{eff}\approx25$ is the effective sample size per LOESS fit.
This gives fixed noise variances in the range 0.0003--0.007
(Table~\ref{tab:hyperparams}).  Unlike a freely estimated
observation-noise term,
$\delta_k^2$ is not optimized; only the signal variance $\alpha_k^2$
and length-scale $\ell_k$ are estimated via marginal likelihood
maximization.

\paragraph{Justification of homoscedastic noise.}
LOESS pre-smoothing absorbs the $\spar$-dependent KDE noise before the
GP sees the data.  What remains---the LOESS estimation uncertainty
$\delta_k^2$---is approximately constant across $\spar$ (varying by a
factor of $<2$ across the training range).  A legacy three-step
Goldberg procedure with input-dependent noise was tested; its large
fixed noise floor swamps the actual LOESS estimation noise
($\sim$0.0003--0.007), forcing the GP to treat real signal as noise.  The
three-way comparison (Table~\ref{tab:gp_comparison}) confirms that the
homoscedastic warped GP outperforms that alternative on
all three variables.

\paragraph{Predictive distribution.}
\begin{align}
  m_k(\spar^*)&=\beta_{0,k}+\beta_{1,k}u^*+
    \mathbf{k}_*^\top(\alpha_k^2\mathbf{K}_{\ell_k}+\delta_k^2\mathbf{I})^{-1}
    \mathbf{r}_k,
  \label{eq:gp_mean}\\
  v_k(\spar^*)&=\alpha_k^2
    -\mathbf{k}_*^\top(\alpha_k^2\mathbf{K}_{\ell_k}+\delta_k^2\mathbf{I})^{-1}
    \mathbf{k}_*,
  \label{eq:gp_var}
\end{align}
where $[\mathbf{k}_*]_j=\alpha_k^2\calK_{\ell_k}(u^*,u_j)$,
$\mathbf{r}_k$ is the residual vector after subtracting the fitted
linear mean, and
$u^*=u(\spar^*)$ is the \ECDF-warped test input obtained by applying
Eq.~\eqref{eq:ecdf} to the held-out value~$\spar^*$.  The variance
$v_k(\spar^*)$ is the posterior variance for the latent smooth score
function, before restoring the residual score variance.

\subsection{Predictive variance via noise decomposition}
\label{sec:noise_decomp}

\paragraph{Motivation.}
The GP epistemic variance $v_k(\spar^*)$ alone captures only
uncertainty in the denoised mean score curve.  The LOESS pre-smoother
suppresses residual score variability before GP fitting; without
correction the emulator predicts a plausible conditional mean density
but cannot produce useful interval estimates.

\paragraph{Procedure.}
The full predictive variance for a new score observation at $\spar^*$
combines GP interpolation uncertainty with the LOESS residual variance:
\begin{equation}
  \sigma^2_{\mathrm{pred},k}(\spar^*) =
  v_k(\spar^*)+\sigma^2_{\mathrm{resid},k},
  \label{eq:pred_var}
\end{equation}
where $v_k(\spar^*)$ is the GP posterior variance
(Equation~\ref{eq:gp_var}) and
$\sigma^2_{\mathrm{resid},k}=\mathrm{Var}[\xi_k-\tilde{\xi}_k]$ is the
per-component residual variance from Equation~\ref{eq:residuals}.
Posterior samples are drawn from
$\mathcal{N}(m_k(\spar^*),\sigma^2_{\mathrm{pred},k}(\spar^*))$ and
back-transformed via Equations~\eqref{eq:mfpca} and~\eqref{eq:sqrt_inv}.
The first term captures uncertainty about where the smooth trend lies
between training points; the second restores the KDE sampling
variability that LOESS removed before GP fitting.  No additional
calibration factor is applied.

\paragraph{Limitations.}
This decomposition rests on two assumptions:
(a)~the LOESS residual variance $\sigma^2_{\mathrm{resid},k}$ is
approximately constant across $\spar$;
(b)~the GP interpolation error and KDE sampling noise are approximately
independent.
Both are simplifications.  The per-variable density-space coverage
---94.1\,\% ($X$), 84.9\,\% ($Y$), and 94.0\,\% ($Z$) at 95\,\%
nominal level---indicates that the decomposition is adequate for $X$
and $Z$ but slightly underestimates uncertainty for $Y$, the variable
with the lowest dynamic range.  The slight under-coverage for $Y$ may
reflect mild $\spar$-dependence of the KDE noise variance that the
homoscedastic model does not capture.

\subsection{Evaluation metrics}
\label{sec:metrics}

Metrics are means over 25 held-out test $\spar$ values.
Formal significance tests (Section~\ref{sec:stat_conf}) and bootstrap
confidence intervals (Table~\ref{tab:stat_tests}) should be consulted
before interpreting directional differences as reliable.

\paragraph{Pre-specified validation criteria.}
We consider the emulator validated if:
\begin{enumerate}
  \item[(a)] leading-PC smoothed-score prediction is accurate
             (positive held-out $R^2$ for PCs~1--8);
  \item[(b)] density-space \iCRPS ratio $<1.0$ for at least two of
             three observables at the 5\,\% significance level;
  \item[(c)] density-space 95\,\% credible-band coverage falls within
             $[85\%,\;100\%]$.
\end{enumerate}

\paragraph{\CRPS in PC space \citep{gneiting2007}.}
\begin{equation}
  \CRPS_k(\spar^*)=
    (y_k-m_k)[2\Phi(z_k)-1]
    +s_k^*[2\phi(z_k)-\pi^{-1/2}],
  \label{eq:crps}
\end{equation}
where $z_k=(y_k-m_k)/s_k^*$, $y_k$ is the empirical mean score,
$s_k^*=\sqrt{v_k(\spar^*)}$ is the GP posterior SD before restoring
the residual variance,
$\Phi$ is the standard normal CDF, and $\phi$ the standard normal PDF.
This metric evaluates the \GP posterior without the restored residual
variance in latent score space.  Because the GP is fit on
LOESS-smoothed scores, $v_k$ alone understates the total predictive
uncertainty; the PC-space CRPS skill score is therefore treated as a
diagnostic rather than a validation criterion.  The primary validation
metrics are the density-space \iCRPS and coverage, which include both
variance components.

\paragraph{\iCRPS --- primary density-space metric \citep{gneiting2007}.}
For the $S=100$ posterior density ensemble $\{f^{(s)}\}$:
\begin{equation}
  \iCRPS(\spar^*)=
    \tfrac{1}{S}\textstyle\sum_s\int|f^{(s)}-f^*|\,\mathrm{d}t
    -\tfrac{1}{2S^2}\textstyle\sum_{s,s'}\int|f^{(s)}-f^{(s')}|
    \,\mathrm{d}t,
  \label{eq:icrps}
\end{equation}
where integrals are approximated via the trapezoidal rule on the
200-point grid.  \iCRPS is a proper scoring rule that evaluates the
full posterior ensemble in the original density space, directly
measuring the emulator's ability to produce calibrated predictive
distributions.  It is the primary density-space validation metric in
this paper because it assesses the full predictive distribution in the
same $L^1$ functional space as the target density, rather than only the
posterior mean.  Density-space point-prediction metrics are therefore
not used for primary emulator validation; reconstruction \ISE is
retained separately only to characterise the basis representation floor
in Table~\ref{tab:mfpca_recon}.

The Energy Score \citep{gneiting2007} was also computed and yields
consistent conclusions; we report \iCRPS as the primary proper scoring
rule.

\paragraph{Credible-band coverage.}
Credible-band coverage at a given test $\spar^*$ is the fraction of the
200 grid points $t$ where $f^*(t)$ falls within the pointwise
$[2.5,97.5]$ quantile envelope of the ensemble $\{f^{(s)}\}$; the
reported coverage is the mean of this fraction over the 25 test $\spar$
values.

\subsection{Naive and LOESS baselines}
\label{sec:naive}

\paragraph{Training-mean naive.}
The simplest conceivable emulator: predict the same density for every
test value of $\spar$, ignoring the input entirely.  Specifically,
it predicts the training-mean density:
\begin{equation}
  \hat{f}_\mathrm{naive}(t)=\frac{1}{99}\sum_j\hat{f}(t|\theta_j).
  \label{eq:naive_mean}
\end{equation}
For the ensemble-based metric (\iCRPS), scores are sampled
from $\mathcal{N}(\bar{\xi}_k,s_{\mathrm{btw},k}^2)$ where
$\bar{\xi}_k$ is the training mean score and $s_{\mathrm{btw},k}^2$
is the between-$\spar$ score variance ($s_{\mathrm{btw},k}^2$ is the
variance of the training mean scores $\{\bar{\xi}_k(\theta_j)\}_{j=1}^{99}$
across the 99 training values; this is distinct from the per-component
residual variance $\sigma^2_{\mathrm{resid},k}$ defined from
Equation~\eqref{eq:residuals}).  This gives the naive
predictor a realistic spread that captures population-level variation;
a zero-spread version would be trivially dominated on proper scoring
rules.  Any emulator that fails to beat this baseline on a given
metric is contributing no information beyond the marginal training
distribution at that metric---a meaningful benchmark precisely because
it sets the lowest acceptable bar.

\paragraph{Score-space LOESS baseline.}
A score-space LOESS regression (span 0.50, degree 1) is fitted to
each PC score vs.\ $\spar$ independently and used to predict mean
scores at test $\spar$ values; densities are reconstructed from these
predictions.  This baseline is deterministic and provides no
uncertainty quantification.  It directly tests whether the GP
machinery adds value over simple nonparametric interpolation for
latent mean prediction.  Because it produces no ensemble, it is not
included in the primary density-space \iCRPS tables.

\paragraph{Nearest-neighbour density baseline.}
For each test $\spar^*$, the $k$ nearest training $\spar$ values
(by Euclidean distance) are identified; their \KDE density estimates
define a simple local density ensemble via bootstrap resampling of the
$k$ neighbours.  This $\spar$-aware conditional density estimator
requires no GP fitting; results for
$k\in\{1,3,5\}$ appear in Table~\ref{tab:nn_comparator}.

\paragraph{Ablated GP variants.}
Two alternative GP configurations are used for ablation
(Table~\ref{tab:gp_comparison}): (i)~a homoscedastic GP without \ECDF
warping (same fixed LOESS-SE$^2$ observation-noise logic), and
(ii)~a heteroscedastic GP without warping (Goldberg 3-step procedure,
Mat\'{e}rn-$\frac{3}{2}$).  Together with the paper configuration
(homoscedastic, \ECDF warping, fixed LOESS-SE$^2$ observation noise
and prediction-time residual variance), these isolate
the marginal contributions of the noise model and input transformation.

Throughout, we report \emph{metric ratios}: for two methods $A$ and
$B$, the ratio $r(A/B) = \overline{m}(A)/\overline{m}(B)$, where
$\overline{m}(\cdot)$ is the mean metric value over 25 test $\spar$
values and lower is better.  A ratio $r < 1$ indicates $A$ outperforms
$B$; $r = 1$ indicates parity.  For binned-mean $R^2$, higher is better,
so ratios above 1 indicate improvement.

\begin{table}[htbp]
  \centering
  \caption{Fixed pipeline hyperparameters (full list).}
  \label{tab:hyperparams}
  \begin{tabular}{lll}
    \toprule
    Parameter & Value & Primary effect \\
    \midrule
    KDE grid points            & 200   & Integral approximation accuracy \\
    KDE B-spline basis         & 80 knots & KDE fidelity \\
    Density floor $\varepsilon$ & $5\times10^{-3}$ & Tail behaviour \\
    MFPCA components $K$       & 30    & Representation accuracy \\
    Eigenfunction basis        & 50 knots & Eigenfunction fidelity \\
    LOESS span (scores)        & 0.25  & Signal/noise split \\
    GP $\ell_\mathrm{min}$     & 0.08  & Over-fitting protection \\
    GP obs.\ noise $\delta_k^2$ & $\sigma^2_{\mathrm{resid},k}/n_\mathrm{eff}$ (fixed) & LOESS estimation uncertainty \\
    Restored variance $\sigma^2_{\mathrm{resid},k}$ & $\mathrm{Var}[\xi_k-\tilde{\xi}_k]$ & KDE sampling noise at prediction \\
    GP mean function           & Linear ($\beta_0+\beta_1 u$) & Handles weak trends \\
    Posterior samples $S$      & 100   & Monte Carlo accuracy \\
    \bottomrule
  \end{tabular}
\end{table}

%% ================================================================
\section{Results}
\label{sec:results}

\subsection{MFPCA representation}
\label{sec:mfpca_quality}

\begin{table}[htbp]
  \centering
  \caption{MFPCA variance explained.  Thirty components account
           for 99.98\,\% of total functional variance.}
  \label{tab:mfpca_var}
  \begin{tabular}{lrr}
    \toprule
    Component & Eigenvalue & Cumul.\ var.\ (\%) \\
    \midrule
    PC\,1  & 2.200 & 21.59 \\
    PC\,2  & 0.682 & 28.28 \\
    PC\,3  & 0.578 & 33.95 \\
    PC\,4  & 0.544 & 39.29 \\
    PC\,5  & 0.476 & 43.96 \\
    PC\,6  & 0.468 & 48.55 \\
    PC\,7  & 0.426 & 52.74 \\
    PC\,8  & 0.389 & 56.57 \\
    PC\,9  & 0.374 & 60.23 \\
    PC\,10 & 0.314 & 63.32 \\
    \midrule
    PC\,20 & 0.174 & 83.63 \\
    PC\,25 & 0.132 & 93.17 \\
    PC\,30 & 0.098 & 99.98 \\
    \bottomrule
  \end{tabular}
\end{table}

\begin{table}[htbp]
  \centering
  \caption{\MFPCA reconstruction \ISE using \emph{exact} unpredicted
           scores (no GP involvement).  These values measure the
           reconstruction accuracy of the 50-knot B-spline basis,
           independently of emulation.}
  \label{tab:mfpca_recon}
  \begin{tabular}{lcc}
    \toprule
    Variable & Mean \ISE & Distributional character \\
    \midrule
    $X$ (Log period ratio)   & 1.308 & Multi-modal, sharp resonance spikes \\
    $Y$ (Log duration ratio) & 0.113 & Smooth, approximately unimodal \\
    $Z$ (Log Hill spacing)   & 0.714 & Moderate resonance structure \\
    \bottomrule
  \end{tabular}
\end{table}

Table~\ref{tab:mfpca_var} shows the variance decomposition.
Table~\ref{tab:mfpca_recon} quantifies the representation floor using
exact, unpredicted scores; this \ISE characterises the basis
independently of emulation and is not used as a primary validation
metric.  The $Y$ distribution (smooth, unimodal) is
reconstructed with mean \ISE\ of 0.113.  The $X$ and $Z$ distributions
contain sharp resonance spikes not well-resolved within the 50-knot
B-spline basis (mean \ISE\ 1.308 and 0.714 respectively).  The
square-root transform improves the representation floor relative to
the CLR-based draft values (1.682 for $X$, 0.114 for $Y$, and 1.147
for $Z$).  These values set the representation floor \emph{for the
current basis choice}; they are not irreducible.  The centre column of
Figure~\ref{fig:reconstruction} illustrates this floor at a
representative test $\spar$.

\subsection{Latent-space GP performance}
\label{sec:gp_results}

\begin{table}[htbp]
  \centering
  \caption{GP emulator diagnostics in latent \MFPCA space (25
           held-out test $\spar$ values).  The CRPS diagnostic uses
           the GP posterior variance before restoring residual variance,
           so density-space metrics remain primary.}
  \label{tab:gp_perf}
  \begin{tabular}{lc}
    \toprule
    Metric & Value \\
    \midrule
    Mean training RMSE                  & 0.0133 \\
    PC1--8 held-out $R^2$ range         & 0.96--0.99 \\
    Binned-mean $R^2$                   & 0.946 \\
    PC-space \CRPS skill score          & $-0.11$ \\
    \bottomrule
  \end{tabular}
\end{table}

Per-component test $R^2$ values range from 0.96 to 0.99 for
PCs~1--8, confirming that the GP accurately predicts the
LOESS-smoothed score curves.  The PC-space \CRPS skill score is
$-0.11$; this reflects the narrow GP posterior variance $v_k$ relative
to the naive baseline spread, not poor mean prediction.  Because the
primary validation metrics (\iCRPS and coverage) operate in density
space where the full predictive variance
$v_k+\sigma^2_{\mathrm{resid},k}$ is used, the latent \CRPS is a
diagnostic rather than a validation criterion.

\subsection{Noise-decomposition coverage}
\label{sec:noise_coverage}

\begin{table}[htbp]
  \centering
  \caption{Density-space credible-band coverage at 95\,\% nominal
           level under the noise-decomposition predictive variance.}
  \label{tab:coverage}
  \begin{tabular}{lccc}
    \toprule
    Metric & $X$ & $Y$ & $Z$ \\
    \midrule
    Density-space 95\,\% coverage & 94.1\,\% & 84.9\,\% & 94.0\,\% \\
    Calibration factor $c$        & 1.00 & 1.00 & 1.00 \\
    \bottomrule
  \end{tabular}
\end{table}

No post-hoc calibration factor is applied ($c=1.0$).  The
noise-decomposition predictive variance (Equation~\ref{eq:pred_var})
yields per-variable coverage of 94.1\,\% ($X$) and 94.0\,\% ($Z$),
near the 95\,\% nominal level.  $Y$ coverage (84.9\,\%) falls just
below the 85\,\% pre-specified threshold, reflecting the variable's
lower dynamic range: with smoother, lower-peaked densities, small
score-space under-dispersion translates to a larger coverage deficit
than for the spiked $X$ and $Z$ distributions.  A per-variable
calibration factor $c_Y\approx1.15$ would restore $Y$ coverage to
95\,\%, but we report the un-tuned result to preserve out-of-sample
validity.

\subsection{Density-space evaluation}
\label{sec:density_eval}

\begin{table}[htbp]
  \centering
  \caption{Density-space evaluation on 25 held-out test $\spar$
           values.  GP columns use the noise-decomposition predictive
           variance (Equation~\ref{eq:pred_var}).  Ratios $<1$ indicate
           GP improvement.  \textbf{Bold}: ratio $<1$; $^*$: $p<0.05$,
           one-sided paired Wilcoxon, $n=25$.}
  \label{tab:density_eval}
  \begin{tabular}{llcccc}
    \toprule
    Metric & Var & GP & Naive & GP/Naive & Verdict \\
    \midrule
    \iCRPS
      & $X$ & 0.837 & 1.003 & \textbf{0.835}\textsuperscript{*} & GP wins \\
      & $Y$ & 0.304 & 0.399 & \textbf{0.761}\textsuperscript{*} & GP wins \\
      & $Z$ & 0.763 & 0.940 & \textbf{0.812}\textsuperscript{*} & GP wins \\[3pt]
    \ISE
      & $X$ & 3.764 & 4.138 & \textbf{0.910} & GP better \\
      & $Y$ & 0.281 & 0.336 & \textbf{0.836} & GP better \\
      & $Z$ & 2.774 & 3.128 & \textbf{0.887} & GP better \\
    \bottomrule
  \end{tabular}
\end{table}

\paragraph{Proper scoring rules: GP wins on all three variables.}
On \iCRPS, the emulator outperforms the naive baseline for all three
observables, with improvements of 17--24\,\% and bootstrap 95\,\% CIs
entirely below 1.0.  This represents a substantial improvement over
the CLR-based pipeline (Section~\ref{sec:ablation_transform}), where
$Y$ showed no significant improvement.

\paragraph{Point prediction: GP modestly improves all variables.}
On \ISE, the GP improves over the naive baseline by 9\,\% ($X$),
16\,\% ($Y$), and 11\,\% ($Z$).  The GP's primary value is in
probabilistic uncertainty quantification (\iCRPS improvement roughly
double the \ISE improvement), not point estimation.

\subsection{Error decomposition}
\label{sec:decomposition}

\begin{table}[htbp]
  \centering
  \caption{ISE decomposition at held-out test $\spar$ values.
           The MFPCA floor uses exact (unpredicted) scores; the GP row
           uses GP-predicted scores.  The difference (GP uplift)
           isolates the GP's contribution to density-space error.}
  \label{tab:decomposition}
  \begin{tabular}{lccc}
    \toprule
    Stage & $X$ ISE & $Y$ ISE & $Z$ ISE \\
    \midrule
    MFPCA floor (true scores)  & 1.308 & 0.113 & 0.714 \\
    GP predicted scores        & 3.764 & 0.281 & 2.774 \\
    Naive baseline             & 4.138 & 0.336 & 3.128 \\
    \midrule
    GP uplift over floor       & +2.456 & +0.168 & +2.060 \\
    GP uplift as \% of total   & 65\,\% & 60\,\% & 74\,\% \\
    \bottomrule
  \end{tabular}
\end{table}

GP score prediction error dominates the remaining density-space error
(60--74\,\% of total ISE), while the MFPCA basis representation
accounts for 26--40\,\%.  This confirms that the basis is adequate for
the current application; future improvement should target GP prediction
quality (more training data, richer kernels) rather than the density
representation.

\subsection{Ablation: density transform}
\label{sec:ablation_transform}

\begin{table}[htbp]
  \centering
  \caption{Effect of density transform on key metrics.  Both
           configurations use MFPCA ($K=30$), no peak alignment,
           and homoscedastic GP with LOESS-SE$^2$ observation noise.
           CLR values come from the previous post-hoc-calibration
           configuration; $\sqrt{f}$ values use the Apr29
           noise-decomposition pipeline.}
  \label{tab:ablation_transform}
  \begin{tabular}{lcc}
    \toprule
    Metric & CLR & $\sqrt{f}$ (Hellinger) \\
    \midrule
    \iCRPS ratio $X$ (GP/Naive) & 0.768 & \textbf{0.835}\textsuperscript{*} \\
    \iCRPS ratio $Y$ (GP/Naive) & 0.983 (n.s.) & \textbf{0.761}\textsuperscript{*} \\
    \iCRPS ratio $Z$ (GP/Naive) & 0.832 & \textbf{0.812}\textsuperscript{*} \\
    Variables with $p<0.05$     & 2 of 3 & \textbf{3 of 3} \\
    Density 95\,\% coverage     & 61.1\,\% & \textbf{94.1/84.9/94.0\,\%} \\
    Calibration factor $c$      & tuned & \textbf{1.00} \\
    GP uplift $X$               & +2.89 & \textbf{+2.46} \\
    \bottomrule
  \end{tabular}
\end{table}

The square-root transform substantially improves calibration and makes
all three \iCRPS comparisons statistically significant.  The most
dramatic effect is on density-space coverage (61\,\% $\to$
near-nominal coverage for $X$ and $Z$), confirming that the CLR's
exponential inverse was the primary driver of miscalibration.  $Y$
(log duration ratio) becomes statistically significant under the
$\sqrt{f}$ transform, reflecting reduced error amplification for this
variable whose densities have lower dynamic range.  The \iCRPS
\emph{ratio} for $X$ appears slightly worse (0.835 vs 0.768) because
the naive baseline also improves without CLR amplification; the
\emph{absolute} GP \iCRPS remains better (0.837 vs 0.847 for $X$).

%% ================================================================
\section{Discussion}
\label{sec:discussion}

\subsection{Interpreting the results}
\label{sec:evidence}

The final pipeline achieves near-calibrated density-space uncertainty:
per-variable coverage is 94.1\,\% ($X$), 84.9\,\% ($Y$), and
94.0\,\% ($Z$) at 95\,\% nominal level, without post-hoc calibration
(Table~\ref{tab:coverage}).  For the $X$ and $Z$ observables---which
exhibit sharp resonance peaks and high dynamic range---coverage is
within 1\,\% of the nominal level.  The $Y$ observable (smooth, low
dynamic range) shows slight under-coverage, potentially addressable by
modelling $\spar$-dependent noise variance.  The absence of a tuned
calibration factor means these coverage values are genuine
out-of-sample assessments.

The GP's primary value is near-calibrated probabilistic prediction.
On \iCRPS, the emulator improves over the naive baseline by 17--24\,\%
for all three variables; on \ISE, posterior-mean improvements are more
modest at 9--16\,\% (Table~\ref{tab:density_eval}).  This asymmetry is
appropriate for an emulator designed for MCMC: proper scoring rules
reward accurate predictive distributions, while point-prediction
metrics ignore whether uncertainty is usable.

The density transform is the most consequential architectural choice.
Replacing CLR with the square-root density transform improves
density-space coverage from 61\,\% to near-nominal values for $X$ and
$Z$ without post-hoc calibration and makes the $Y$ comparison
statistically significant (Table~\ref{tab:ablation_transform}).
The improvement is larger than the differences among GP noise-model
variants (Table~\ref{tab:gp_comparison}).

Removing curve registration also simplified the pipeline and eliminated
a dominant error source.  The final MFPCA basis absorbs peak movement
through its eigenfunctions rather than requiring separate geometry
parameters that must themselves be emulated at test $\spar$ values.

\subsection{Representation floor: current basis, not fundamental limit}
\label{sec:discussion_floor}

The representation floor is specific to the 50-knot B-spline basis and
$K=30$ components.  The exact-score reconstruction \ISE
(Table~\ref{tab:mfpca_recon}: 1.308 for $X$, 0.113 for $Y$, and
0.714 for $Z$), computed without any GP involvement, establishes that
this floor is a property of the basis, not of the emulator.  The error
decomposition (Table~\ref{tab:decomposition}) shows that the bottleneck
has shifted: GP score prediction accounts for 60--74\,\% of total
density-space \ISE, while the basis representation accounts for
26--40\,\%.  Future accuracy gains should therefore target score
prediction quality and the residual-noise model first.

\subsection{Sensitivity to pipeline choices}
\label{sec:sensitivity}

The transform ablation (Table~\ref{tab:ablation_transform}) dominates
the sensitivity story.  With the final square-root representation fixed,
the three-way GP comparison in Table~\ref{tab:gp_comparison} shows that
the homoscedastic warped GP is competitive with, and usually better
than, the no-warp and heteroscedastic alternatives.  This supports the
two-stage design: LOESS removes high-frequency KDE noise first, and the
GP emulates the denoised score curve with the LOESS estimation
uncertainty as fixed known observation noise.  The residual variance
stripped by LOESS is restored at prediction time, giving a principled
uncertainty decomposition that requires no post-hoc calibration.

\subsection{Baseline adequacy}
\label{sec:benchmarks}

Three baselines delimit where the GP adds value.

\begin{table}[htbp]
  \centering
  \caption{Nearest-neighbour density comparison ($k=1,3,5$ neighbours)
           vs GP, mean over 25 test $\spar$ values, reported via
           \iCRPS (lower is better).  The NN comparator is converted
           into a simple local ensemble by bootstrap resampling the
           $k$ nearest training densities; GP/NN-k3$<1$ means the full
           GP beats the 3-nearest-neighbour density ensemble.  NN
           columns are unchanged by the GP noise-model update.}
  \label{tab:nn_comparator}
  \setlength{\tabcolsep}{4pt}
  \begin{tabular}{lccccc}
    \toprule
    Metric & NN-$k$1 & NN-$k$3 & NN-$k$5 & GP & GP/NN-$k$3 \\
    \midrule
    \iCRPS\ $X$      & 0.820 & 1.061 & 0.975 & 0.837 & 0.789 \\
    \iCRPS\ $Y$      & 0.323 & 0.404 & 0.366 & 0.304 & 0.752 \\
    \iCRPS\ $Z$      & 0.807 & 1.036 & 0.934 & 0.763 & 0.737 \\
    \bottomrule
  \end{tabular}
\end{table}

\begin{table}[htbp]
  \centering
  \caption{Three-way \GP comparison isolating the contribution of
           noise model choice and \ECDF input warping.
           GP-HO-NW: homoscedastic, no warping.
           GP-HE-NW: heteroscedastic (Goldberg 3-step), no warping.
           GP-HO-W: homoscedastic with \ECDF warping, fixed LOESS-SE$^2$
           observation noise, and prediction-time residual variance
           (paper configuration).
           \iCRPS: lower is better; binned-mean $R^2$: higher is better.
           Ratios $<1$ indicate improvement of GP-HO-W over the comparator.}
  \label{tab:gp_comparison}
  \setlength{\tabcolsep}{4pt}
  \begin{tabular}{lccccc}
    \toprule
    Metric & GP-HO-NW & GP-HE-NW & GP-HO-W & HO-W/HO-NW & HO-W/HE-NW \\
    \midrule
    \iCRPS\ $X$          & 1.207 & 1.173 & 0.837 & 0.693 & 0.714 \\
    \iCRPS\ $Y$          & 0.426 & 0.426 & 0.304 & 0.712 & 0.713 \\
    \iCRPS\ $Z$          & 1.104 & 1.083 & 0.763 & 0.691 & 0.705 \\
    Binned-mean $R^2$    & 0.988 & 0.439 & 0.946 & 0.958 & --- \\
    \bottomrule
  \end{tabular}
\end{table}

Against the $\spar$-aware 3-nearest-neighbour estimator, the GP
matches or outperforms on \iCRPS for all three variables
(GP/NN-k3: 0.737--0.789; Table~\ref{tab:nn_comparator}), establishing
that the model adds value beyond simple conditional density averaging
while also providing smooth interpolation and uncertainty bands.

The three-way GP comparison (Table~\ref{tab:gp_comparison}) isolates
architectural contributions.  GP-HO-NW removes input warping, while
GP-HE-NW replaces the final homoscedastic observation-noise model with a legacy
heteroscedastic model.  GP-HO-W (the paper configuration) achieves the
lowest \iCRPS on all three variables, supporting the final choice of
\ECDF warping plus the noise-decomposition predictive variance.

Comparing univariate FPCA per density against \MFPCA would isolate
the value of joint modelling; given that the primary remaining error
now traces to score prediction, this comparison is deferred to future
work.

\subsection{Significance testing}
\label{sec:stat_conf}

\begin{table}[htbp]
  \centering
  \caption{Significance tests for \iCRPS ($n=25$ test pairs).
           Wilcoxon: one-sided $H_1$: GP $<$ Naive
           (lower \iCRPS = better).
           Bootstrap CI: 95\,\% interval on GP/Naive ratio
           (5000 resamples); CI entirely below 1.0 confirms GP wins.}
  \label{tab:stat_tests}
  \begin{tabular}{lccc}
    \toprule
    Variable & $p$-value & Sig. & 95\,\% CI (ratio) \\
    \midrule
    $X$ (Log period ratio)   & $<$0.0001 & $*$ & [0.810, 0.859] \\
    $Y$ (Log duration ratio) & $<$0.0001 & $*$ & [0.683, 0.833] \\
    $Z$ (Log Hill spacing)   & $<$0.0001 & $*$ & [0.779, 0.848] \\
    \bottomrule
  \end{tabular}

  \smallskip
  {\small $*$ $p<0.05$.  All three bootstrap CIs lie entirely below
  1.0, confirming the GP outperforms the naive baseline.}
\end{table}

All three \iCRPS comparisons ($X$, $Y$, $Z$) are statistically
significant at the 5\,\% level; bootstrap 95\,\% confidence intervals
on the GP/Naive ratio lie entirely below~1.0 for all three
(Table~\ref{tab:stat_tests}).  Point-prediction improvements are
modest but consistent, while the probabilistic improvements are
statistically decisive.

\subsection{Alternative density representations}
\label{sec:alternatives}

For $X$ (and to a lesser extent $Z$), the following approaches may
reduce basis-limitation error: location-scale parametrisation of
resonance spikes \citep{ma2011}; amplitude-phase separation before
\MFPCA \citep{ramsay2005, kurtek2011}; quantile-function
representation; or Gaussian mixture emulation.  These alternatives are
secondary to the immediate bottleneck identified in
Table~\ref{tab:decomposition}: score prediction error now accounts for
most of the remaining density-space \ISE.

%% ================================================================
\section{Conclusions}
\label{sec:conclusions}

A Gaussian process emulation workflow for multivariate conditional
density functions from planetary system simulations yields three
principal findings.

\emph{The pipeline produces near-calibrated probabilistic density
predictions.}  Per-variable density-space 95\,\% credible-band
coverage reaches 94.1\,\% ($X$) and 94.0\,\% ($Z$) without post-hoc
calibration, confirming that the noise-decomposition predictive
variance provides reliable uncertainty quantification for downstream
Bayesian inverse inference.  $Y$ coverage (84.9\,\%) falls slightly
below the pre-specified 85\,\% threshold, attributable to the variable's
lower dynamic range; this is noted as a limitation.  The \iCRPS proper
scoring rule shows statistically significant improvement over the naive
baseline for all three observables (17--24\,\% improvement, all
$p<0.05$).

\emph{The square-root density transform is the critical architectural
choice.}  Replacing the standard CLR transform with the Hellinger
($\sqrt{f}$) transform eliminates the exponential error amplification
in the CLR inverse, improving density-space coverage from 61\,\% to
near-nominal values for $X$ and $Z$ without post-hoc calibration and
enabling $Y$ (log duration ratio) to achieve statistical significance.
This insight may be broadly applicable to other
functional data emulation problems involving density-valued outputs.

\emph{GP score prediction is the primary remaining error source.}
An ISE decomposition shows that GP prediction error accounts for
60--74\,\% of total density-space ISE, with the MFPCA basis
representation contributing 26--40\,\%.  Future improvement should
target GP prediction quality (more training simulations, richer
kernels) and mild $\spar$-dependence in the residual-noise variance
rather than the density representation.

The emulator is ready for deployment in Bayesian inverse inference of
planet formation conditions from Kepler observations; this application
is the subject of a companion paper.

%% ================================================================
\section*{Acknowledgements}
The author thanks the developers and maintainers of the open-source
\textsf{R} packages used in this analysis.  No external funding was
reported for this work.

\section*{Data Availability}
The simulation data and \textsf{R} analysis code sufficient to
reproduce all tables and figures in this paper will be deposited in a
public repository at acceptance.  The analysis was conducted in
\textsf{R}~4.3.2; key packages are \texttt{fda}~6.1.4
\citep{ramsay2005}, \texttt{MFPCA}~1.3 \citep{happ2018},
\texttt{DiceKriging}~1.6.0 \citep{roustant2012},
\texttt{scoringRules}~1.1.3 \citep{jordan2019}, and
\texttt{ggplot2}~3.4.4.

%% ================================================================
\begin{thebibliography}{99}

\bibitem[Abramson(1982)]{abramson1982}
Abramson, I.~S.\ 1982, \textit{Ann.\ Stat.}, 10, 1217.

\bibitem[Benz et al.(2014)]{benz2014}
Benz, W., Ida, S., Alibert, Y., Lin, D., \& Mordasini, C.\ 2014,
in \textit{Protostars and Planets VI}, eds.\ H.\ Beuther et al.,
University of Arizona Press, p.\ 691.

\bibitem[Bower et al.(2010)]{bower2010}
Bower, R.~G.\ et al.\ 2010, \textit{MNRAS}, 407, 2017.

\bibitem[Egozcue et al.(2003)]{egozcue2003}
Egozcue, J.~J., Pawlowsky-Glahn, V., Mateu-Figueras, G., \&
Barcel\'{o}-Vidal, C.\ 2003, \textit{Math.\ Geol.}, 35, 279.

\bibitem[Emsenhuber et al.(2021)]{emsenhuber2021}
Emsenhuber, A.\ et al.\ 2021, \textit{A\&A}, 656, A69.

\bibitem[Fabrycky et al.(2014)]{fabrycky2014}
Fabrycky, D.~C.\ et al.\ 2014, \textit{ApJ}, 790, 146.

\bibitem[Fricker et al.(2013)]{fricker2013}
Fricker, T.~E., Oakley, J.~E., \& Urban, N.~M.\ 2013,
\textit{Technometrics}, 55, 174.

\bibitem[Gilbert \& Fabrycky(2020)]{gilbert2020}
Gilbert, G.~J.\ \& Fabrycky, D.~C.\ 2020, \textit{AJ}, 159, 281.

\bibitem[Gneiting \& Raftery(2007)]{gneiting2007}
Gneiting, T.\ \& Raftery, A.~E.\ 2007,
\textit{J.\ Am.\ Stat.\ Assoc.}, 102, 359.

\bibitem[Goldberg et al.(1998)]{goldberg1998}
Goldberg, P.~W., Williams, C.~K.~I., \& Bishop, C.~M.\ 1998,
in \textit{Advances in Neural Information Processing Systems 10},
eds.\ M.~I.\ Jordan, M.~J.\ Kearns \& S.~A.\ Solla,
MIT Press, pp.\ 493--499.

\bibitem[Happ \& Greven(2018)]{happ2018}
Happ, C.\ \& Greven, S.\ 2018, \textit{Biometrics}, 74, 237.

\bibitem[Hayashi(1981)]{hayashi1981}
Hayashi, C.\ 1981, \textit{Prog.\ Theor.\ Phys.\ Suppl.}, 70, 35.

\bibitem[Higdon et al.(2008)]{higdon2008}
Higdon, D., Gattiker, J., Williams, B., \& Rightley, M.\ 2008,
\textit{J.\ Am.\ Stat.\ Assoc.}, 103, 570.

\bibitem[Ida \& Lin(2004)]{ida2004}
Ida, S.\ \& Lin, D.~N.~C.\ 2004, \textit{ApJ}, 604, 388.

\bibitem[Jordan et al.(2019)]{jordan2019}
Jordan, A., Kr\"uger, F., \& Lerch, S.\ 2019,
\textit{J.\ Stat.\ Softw.}, 90, 1.

\bibitem[Kennedy \& O'Hagan(2001)]{kennedy2001}
Kennedy, M.~C.\ \& O'Hagan, A.\ 2001,
\textit{J.\ R.\ Stat.\ Soc.\ B}, 63, 425.

\bibitem[Kersting et al.(2007)]{kersting2007}
Kersting, K., Plagemann, C., Pfaff, P., \& Burgard, W.\ 2007,
in \textit{Proc.\ 24th ICML}, ACM, pp.\ 393--400.

\bibitem[Kurtek et al.(2011)]{kurtek2011}
Kurtek, S., Srivastava, A., \& Wu, W.\ 2011,
in \textit{Proc.\ CVPR 2011}, IEEE, pp.\ 1995--2002.

\bibitem[Ma \& Zabaras(2011)]{ma2011}
Ma, X.\ \& Zabaras, N.\ 2011, \textit{J.\ Comput.\ Phys.}, 230, 8604.

\bibitem[Mordasini(2018)]{mordasini2018}
Mordasini, C.\ 2018, in \textit{Handbook of Exoplanets},
eds.\ H.~J.\ Deeg \& J.~A.\ Belmonte, Springer, pp.\ 2425--2474.

\bibitem[Mulders et al.(2021)]{mulders2021}
Mulders, G.~D.\ et al.\ 2021, \textit{ApJ}, 920, 66.

\bibitem[Ramsay \& Silverman(2005)]{ramsay2005}
Ramsay, J.~O.\ \& Silverman, B.~W.\ 2005,
\textit{Functional Data Analysis}, 2nd edn., Springer.

\bibitem[Roustant et al.(2012)]{roustant2012}
Roustant, O., Ginsbourger, D., \& Deville, Y.\ 2012,
\textit{J.\ Stat.\ Softw.}, 51, 1.

\bibitem[Rasmussen \& Williams(2006)]{rasmussen2006}
Rasmussen, C.~E.\ \& Williams, C.~K.~I.\ 2006,
\textit{Gaussian Processes for Machine Learning}, MIT Press.

\bibitem[Silverman(1986)]{silverman1986}
Silverman, B.~W.\ 1986,
\textit{Density Estimation for Statistics and Data Analysis},
Chapman \& Hall.

\bibitem[Srivastava et al.(2007)]{srivastava2007}
Srivastava, A., Jermyn, I.~H., \& Joshi, S.~H.\ 2007,
in \textit{Proc.\ IEEE Conf.\ Computer Vision and Pattern Recognition
(CVPR)}, doi:10.1109/CVPR.2007.383188.

\bibitem[van den Boogaart \& Tolosana-Delgado(2014)]{van2014}
van den Boogaart, K.~G.\ \& Tolosana-Delgado, R.\ 2014,
\textit{Analyzing Compositional Data with R}, Springer.

\bibitem[Weiss et al.(2018)]{weiss2018}
Weiss, L.~M.\ et al.\ 2018, \textit{AJ}, 155, 48.

\bibitem[Winn \& Fabrycky(2015)]{winn2015}
Winn, J.~N.\ \& Fabrycky, D.~C.\ 2015, \textit{ARA\&A}, 53, 409.

\bibitem[MacDonald et al.(2020)]{macdonald2020}
MacDonald, M.~G., Dawson, R.~I., Morrison, S.~J., Lee, E.~J., \&
Khandelwal, A.\ 2020, \textit{ApJ}, 891, 20.

\end{thebibliography}

%% ================================================================
\appendix
\section{Figures}
\label{app:figs}

\begin{nolinenumbers}
\begin{figure}[htbp]
  \centering
  \tikzset{
    block/.style  = {rectangle, draw, rounded corners=3pt,
                     fill=blue!8, text width=6.8cm, align=center,
                     minimum height=0.75cm, font=\small},
    arrowstyle/.style = {-{Stealth[length=5pt]}, thick, draw=gray!70},
    sidelab/.style = {font=\scriptsize, text=gray!80, align=left}
  }
  \begin{tikzpicture}[node distance=0.55cm]
    \node[block, fill=orange!15] (input)
      {Input: $\spar$ at held-out value $\spar^*$};
    \node[block, below=of input] (sim)
      {$N$-body simulations at training $\theta_j$ \\
       $\Rightarrow$ system catalogues $(X,Y,Z)$};
    \node[block, below=of sim] (kde)
      {Adaptive \KDE\ on 200-pt grids\\
       $\Rightarrow$ $\hat{f}_{X}(\cdot|\theta_j),\;\hat{f}_{Y},\;\hat{f}_{Z}$};
    \node[block, below=of kde] (sqrt)
      {$\sqrt{f}$ transform\\
       $h_v = \sqrt{\hat{f}_v}$\quad($v\in\{X,Y,Z\}$)};
    \node[block, below=of sqrt] (mfpca)
      {\MFPCA\ ($K=30$ components, 50-knot B-splines)\\
       $\Rightarrow$ scores $\{\xi_k(\theta_j)\}_{k=1}^{30}$};
    \node[block, below=of mfpca] (loess)
      {LOESS pre-smooth (span $0.25$) $\Rightarrow$ $\tilde{\xi}_k$;\\
       residuals $r_k = \xi_k - \tilde{\xi}_k$
       $\Rightarrow$ $\sigma^2_{\mathrm{resid},k}$};
    \node[block, fill=green!8, below=of loess] (gp)
      {Homoscedastic \GP\ per $k$: Mat\'ern-$\tfrac{3}{2}$,
       \ECDF\ warping, fixed LOESS-SE$^2$ observation noise\\
       $\Rightarrow$ posterior mean $m_k(\spar^*)$, variance $v_k(\spar^*)$};
    \node[block, below=of gp] (predvar)
      {Predictive variance:
       $\sigma^2_{\mathrm{pred},k}=v_k(\spar^*)+\sigma^2_{\mathrm{resid},k}$\\
       Draw $\hat{\xi}_k^{(s)}\sim
       \mathcal{N}(m_k,\sigma^2_{\mathrm{pred},k})$,\quad$s=1,\ldots,100$};
    \node[block, below=of predvar] (recon)
      {\MFPCA\ back-reconstruction:\\
       $\hat{h}^{(s)}(\cdot)=\boldsymbol{\mu}+\sum_k\hat{\xi}_k^{(s)}\boldsymbol{\phi}_k$};
    \node[block, fill=orange!15, below=of recon] (output)
      {Inverse $\sqrt{f}$: $f=h^2/\int h^2$\\
       $\{\hat{f}_v^{(s)}(\cdot|\spar^*)\}_{s=1}^{100}$\quad($v\in\{X,Y,Z\}$)};
    \draw[arrowstyle] (input)  -- (sim);
    \draw[arrowstyle] (sim)    -- (kde);
    \draw[arrowstyle] (kde)    -- (sqrt);
    \draw[arrowstyle] (sqrt)   -- (mfpca);
    \draw[arrowstyle] (mfpca)  -- (loess);
    \draw[arrowstyle] (loess)  -- (gp);
    \draw[arrowstyle] (gp)     -- (predvar);
    \draw[arrowstyle] (predvar) -- (recon);
    \draw[arrowstyle] (recon)  -- (output);
  \end{tikzpicture}
  \caption{Pipeline overview.  Simulation outputs are converted to
    densities via adaptive \KDE, mapped to square-root density functions,
    jointly compressed by \MFPCA, emulated in score space by a
    homoscedastic \GP with \ECDF input warping and fixed LOESS-SE$^2$
    observation noise, sampled with the noise-decomposition predictive
    variance (Equation~\ref{eq:pred_var}), and back-transformed by
    squaring and normalising.
    Orange boxes are inputs/outputs; green box is the core emulation step.}
  \label{fig:pipeline}
\end{figure}
\end{nolinenumbers}

\begin{figure}[htbp]
  \centering
  \includegraphics[width=0.90\textwidth]{fig_scores.pdf}
  \caption{\MFPCA\ score series for PCs~1--8 against the simulation
    control parameter $\theta$ (g\,cm$^{-2}$).
    Grey dots: per-$\theta$ mean raw scores $\xi_k(\theta)$;
    blue line: LOESS-smoothed $\tilde{\xi}_k(\theta)$ (span $=0.25$).
    Each panel shows the lag-1 autocorrelation $r_1$ of the smoothed
    series and the residual SD defining
    $\sigma^2_{\mathrm{resid},k}$.
    All components show $r_1>0.91$ ($r_1=0.92$--$0.98$ across PCs~1--8),
    confirming structured $\theta$-dependence throughout.}
  \label{fig:scores}
\end{figure}

\begin{figure}[htbp]
  \centering
  \includegraphics[width=\textwidth]{fig_reconstruction.pdf}
  \caption{Three-stage reconstruction at representative held-out
    $\spar^*=70.4$\,g\,cm$^{-2}$ (near-median \iCRPS for $Y$).
    \textbf{Left:} held-out adaptive \KDE (the emulation target).
    \textbf{Centre:} \MFPCA reconstruction from exact (unpredicted)
    scores, showing the basis representation floor independent of \GP
    quality; \ISE annotated per panel.
    \textbf{Right:} full \GP emulator output with 95\,\% credible
    band (light blue); \iCRPS and pointwise coverage annotated per panel.
    Black dashed line in centre and right columns reproduces the \KDE
    target for visual comparison.}
  \label{fig:reconstruction}
\end{figure}

\begin{figure}[htbp]
  \centering
  \includegraphics[width=0.97\textwidth]{fig_ppd.pdf}
  \caption{Posterior predictive densities at four
    held-out $\spar$ values (columns), for each of
    the three planetary-system observables (rows).
    \textbf{Black}: held-out \KDE; \textbf{red dashed}: training-mean
    naive baseline; \textbf{blue}: \GP\ posterior mean;
    \textbf{shaded}: 95\,\% credible band.  Empirical density-space
    coverage is annotated in each panel (nominal 95\,\%).}
  \label{fig:ppd}
\end{figure}

\end{document}
